\scanpage{33}
\inserted{$\Omega(s_1,s_{\unsure{2|7}})$} défini par le couple de sections $(s_1,s_2)$,
d'idéal $\unsure{\underline{\mathscr I}|\underline{\mathscr J}|\underline{\mathcal J}}(s_1)\cdot\underline{\mathscr I}(s_2)\subset\underline O_F$, en la
section
\[
\underline b_{\Omega(s_1,s_2)}=s_1+s_2\in\Gamma(2F). \tag*{(41)}
\]
En termes d'un $\Omega$ quelconque, on
peut aussi dire que $\underline b_\Omega$ est caractérisé par
la condition suivante : pour tte \deleted{splitting}\inserted{section}
(\unsure{évt locale|soit locale}) $s$ de $F$, considérant la section
$2s$ de $2F$, on a
\[
 b(s)=\underline b(s)-2s \tag*{(42)}
\]
Si $\Omega=\Omega(s_1,s_2)$, cela donne bien
\[
 b(s)=s_1+s_2-2s=(s_1-s)+(s_2-s) \tag*{(43)}
\]

On peut donc dire que la donnée
d'une section $\sigma$ de $\Sigma^2(F)$ --- i.e. d'un
diviseur quadratique $\Omega\subset F$ --- revient à
la donnée successive

a) d'une section $\underline b$ de $2F$ --- d'où
par \deleted{\unsure{b) dé}} image inverse un fibré affine de
rang $1$
\[
 \unsure{\mathbb X|\Sigma|\overline\Sigma}(\underline b)\quad\deleted{\text{du groupe}} \tag*{(44)}
\]
sous le groupe des translations $L^{\otimes2}$ ; et

b) d'une section $\sigma$ de $\mathbb X(\underline b)$.

% archive sheet number: 32

